\section*{Bab \ref{ch:b3:relativistic-dynamics} --- Dinamika Relativistik}

\begin{solution}{exo:b3:relativistic-dynamics:1}
Elektron: $\gamma = 7.09$; $pc = \gamma\beta\,mc^2 = \qty{3.59}{MeV}$,
sehingga $p = \qty{3.59}{MeV}/c$; $E = \qty{3.62}{MeV}$, $E_k =
\qty{3.11}{MeV}$. Proton: $E = 938 + 1000 = \qty{1938}{MeV}$, $\gamma
= 2.07$, $pc = \sqrt{E^2 - (mc^2)^2} = \qty{1696}{MeV}$, $\beta =
pc/E = 0.875$.
\end{solution}

\begin{solution}{exo:b3:relativistic-dynamics:2}
(a) $\num{e-3} \times \num{9e16} = \qty{9e13}{J}$. (b) $\num{6e13}/
\num{9e16} \approx \qty{0.7}{g}$ --- bomnya mengubah kurang dari satu
gram. (c) $L/c^2 = \qty{4.2e9}{kg/s}$; selama \qty{e10}{yr}
($\num{3.2e17}{}\,\unit{s}$): $\num{1.3e27}{}\,\unit{kg}$, sekitar
$0.07\%$ Matahari. (d) $\num{5e4}/\num{9e16} \approx
\qty{0.6}{ng}$ --- nyata, dan selamanya tak tertimbang.
\end{solution}

\begin{solution}{exo:b3:relativistic-dynamics:3}
(a) $p = P/c$ tiap sekon: gaya $F = P/c = \qty{1.7e-11}{N}$ berupa
hentakan balik. (b) Pemantulan melipatduakan pemindahannya: $F = 2\Phi A/c =
2 \times 1400 \times 100/\num{3e8} \approx \qty{0.9}{mN}$. (c) $a
\approx \qty{9e-5}{m/s^2}$; setelah \qty{2.6e6}{s}: $v \approx
\qty{240}{m/s}$ --- lambat pada awalnya, tetapi mesinnya tak pernah
kehabisan bahan bakar. (d) Momentum sinar matahari (bersama angin
surya) mendorong debunya terus-menerus ke luar: ekornya mengalir
menjauhi Matahari, bukan tertinggal di belakang kometnya.
\end{solution}

\begin{solution}{exo:b3:relativistic-dynamics:4}
(a) $m = E_0/c^2$: $\num{0.511e6} \times \num{1.6e-19}/\num{9e16} =
\qty{9.1e-31}{kg}$. (b) $\num{e9} \times \num{1.6e-19}/\num{3e8} =
\qty{5.3e-19}{kg.m/s}$. (c) $\sqrt s = \qty{4}{GeV}$ energi pusat
momentum. (d) Dengan $c = 1$, energi, momentum dan massa berbagi satu
satuan dan tiap rumus melepaskan $c$-nya; untuk kembali ke SI,
masukkanlah kembali pangkat $c$ tunggal yang membetulkan dimensinya.
\end{solution}

\begin{solution}{exo:b3:relativistic-dynamics:5}
(a) $2E_k/mc^2$: $1.96$, $3.9$, $17.6$, $59$ --- mustahil setelah nilai
yang pertama. (b) $1 - (1 + E_k/mc^2)^{-2}$: $0.74$, $0.89$, $0.990$,
$0.9989$. (c) $8.4/0.99946c = \qty{28.0}{ns}$; nilai Newton $7.7c$ akan
terbaca \qty{3.7}{ns}. (d) Kalorimeternya membuktikan bahwa elektronnya
sungguh membawa seluruh \qty{15}{MeV} ke dalam sasarannya: energinya
memang ada di sana, namun lajunya sudah berhenti tumbuh --- usaha kini
membeli momentum dan energi, bukan kecepatan.
\end{solution}

\begin{solution}{exo:b3:relativistic-dynamics:6}
(a) $m_\pi c^2 = \sqrt{p^2c^2 + m_\mu^2c^4} + pc$; kucilkanlah akarnya
lalu kuadratkan: $pc = (m_\pi^2 - m_\mu^2)c^4/2m_\pi c^2 =
(139.6^2 - 105.7^2)/(2 \times 139.6) = \qty{29.8}{MeV}$. (b) $E_\mu =
\sqrt{29.8^2 + 105.7^2} = \qty{109.8}{MeV}$: $E_k = \qty{4.1}{MeV}$,
$\beta = 29.8/109.8 = 0.27$. (c) $E_\nu = pc = \qty{29.8}{MeV}$. (d)
Dengan mendorongnya: $E_{\text{lab}} = \gamma(E^* \pm \beta p^*c)$ dengan $\beta
\approx 1$: antara $50 \times (109.8 - 29.8) = \qty{4.0}{GeV}$ dan
$50 \times (109.8 + 29.8) = \qty{7.0}{GeV}$.
\end{solution}

\begin{solution}{exo:b3:relativistic-dynamics:7}
(a) Untuk dua kuantum tak bermassa, $M^2c^4 = 2E_1E_2(1 - \cos\theta) =
2 \times 200^2 \times \tfrac12$: $M = \qty{200}{MeV}/c^2$. (b)
Tepat $m_{\pi^0}$ --- \omterm{def:b3:relativistic-dynamics:four-momentum}{massa invarian} adalah massa partikelnya, dalam
setiap kerangka. (c) Hitunglah $M$ untuk tiap pasangan foton dalam
peristiwanya: pasangan acak tersebar mulus, sedangkan anak sejati sebuah
partikel menumpuk pada massanya --- itulah tonjolannya. (d) $\theta = 0$:
$M = 0$; berkas cahaya sejajar tidak mempunyai kerangka diam --- ia
bergerak dengan $c$ sebagai satu kesatuan, seperti tiap fotonnya.
\end{solution}

\begin{solution}{exo:b3:relativistic-dynamics:8}
(a) $\gamma = \num{6.8e12}/\num{9.38e8} = 7250$; $c - v \approx
c/2\gamma^2 = \qty{2.9}{m/s}$. (b) $f = v/L \approx \num{3e8}/27000 =
\qty{11.1}{kHz}$: sebelas ribu putaran tiap sekon. (c) $\num{3e14}
\times \qty{6.8}{TeV} = \qty{3.3e8}{J} \approx \qty{78}{kg}$ TNT
--- dalam berkas setipis rambut. (d) $2 \times \num{3.3e8}/\num{9e16} \approx
\qty{7}{\micro g}$: persediaan kerja materi yang akan dibuat.
\end{solution}

\begin{solution}{exo:b3:relativistic-dynamics:9}
(a) \omterm{def:b3:relativistic-dynamics:four-momentum}{Massa invarian} sebuah foton adalah $0$; \omterm{def:b3:relativistic-dynamics:four-momentum}{massa invarian} pasangan
$e^+e^-$ paling sedikit $2m_{\text{e}}c^2$. Invariannya tidak dapat
berubah dalam peluruhan terkucil: terlarang. (Setara dengan itu: tak ada
kerangka yang membuat momentumnya setimbang.) (b)
Berhadapan, $M^2c^4 = 4E\epsilon$; penciptaan pasangannya menuntut $Mc^2 \ge
2m_{\text{e}}c^2$, yakni $E\epsilon \ge (m_{\text{e}}c^2)^2$. (c)
Foton \qty{511}{keV} yang lain; untuk \qty{50}{keV}, pasangan sebesar
$0.511^2/0.05 = \qty{5.2}{MeV}$. (d) Foton TeV memenuhi ambang
$(m_{\text{e}}c^2)^2/E \sim \qty{0.3}{eV}$ pada cahaya bintang biasa:
langit itu sendiri menyerap sinar $\gamma$ TeV pada jarak kosmologis ---
semesta tidak tembus pandang pada semua energi.
\end{solution}

\begin{solution}{exo:b3:relativistic-dynamics:10}
(a) $P_{e'} = P_\gamma + P_e - P_{\gamma'}$; kuadratkanlah kedua ruasnya
(invarian): $m_{\text{e}}^2c^4 = m_{\text{e}}^2c^4 + 2P_\gamma\!
\cdot\!P_e - 2P_{\gamma'}\!\cdot\!P_e - 2P_\gamma\!\cdot\!P_{\gamma'}$.
(b) Dengan $P_\gamma\!\cdot\!P_e = E m_{\text{e}}$,
$P_{\gamma'}\!\cdot\!P_e = E'm_{\text{e}}$ dan $P_\gamma\!\cdot\!
P_{\gamma'} = EE'(1 - \cos\theta)/c^2$:
$m_{\text{e}}c^2(E - E') = EE'(1 - \cos\theta)$, yang dalam panjang
gelombang ($E = hc/\lambda$) merupakan pergeseran yang dinyatakan itu.
(c) $h/m_{\text{e}}c =
\qty{2.43}{pm}$, paling besar \qty{4.9}{pm}:
satu per $10^5$ dari cahaya tampak (tak terlihat), beberapa persen dari
sinar-X \qty{100}{pm} (terukur). (d) Bahwa sebuah foton menumbuk
bagaikan bola bilyar --- membawa momentum $h/\lambda$ yang dipertukarkan
dalam peristiwa \emph{perorangan}, bukan sekadar energi dalam bongkahan.
\end{solution}

\begin{solution}{exo:b3:relativistic-dynamics:11}
(a) $M^2c^4 = m_{\text{p}}^2c^4 + 4E\epsilon$ secara berhadapan (dengan
protonnya ultrarelativistik); ambangnya pada $M = m_\Delta$. (b) $E =
(1232^2 - 938^2)\,\unit{MeV^2}/(4 \times \num{6e-4}\,\unit{eV})
\approx \qty{2.7e20}{eV}$ untuk foton rata-ratanya --- ekor termal
latar gelombang mikro menurunkan pemenggalan efektifnya menjadi
$\sim\qty{6e19}{eV}$. (c) Proton di atas pemenggalannya kehilangan
energi kepada $\Delta$ dalam $\sim\qty{e8}{ly}$: spektrumnya seharusnya
berakhir di dekat \qty{6e19}{eV} bagi sumber yang jauh --- yakni
penekanan GZK yang teramati.
(d) Sumbernya haruslah sekaligus pemercepat luar biasa dan \emph{dekat}
secara kosmis --- di dalam supergugus kita --- dan itulah sebagian
sebabnya asal usulnya masih diburu.
\end{solution}

\begin{solution}{exo:b3:relativistic-dynamics:12}
(a) Energi: $m_{\text{i}}c^2 = \gamma m_{\text{f}}c^2 + E_\ell$;
momentum: $\gamma m_{\text{f}}\beta c = E_\ell/c$. Lenyapkanlah
$E_\ell$: $m_{\text{i}} = \gamma(1 + \beta)m_{\text{f}} =
m_{\text{f}}\sqrt{(1+\beta)/(1-\beta)}$. (b) $\sqrt{19} = 4.4$;
mempercepat \emph{dan} mengerem mengkuadratkannya: $19$. (c) Massa yang
dihabiskan $m_{\text{i}} - m_{\text{f}} = 3.4\,m_{\text{f}}$ haruslah
bahan bakar yang dimusnahkan, separuhnya antimateri: \qty{1.7}{kg}
antimateri tiap kilogram yang diantarkan (sekali jalan, tanpa
pengereman). (d) Dengan produksi sedunia beberapa nanogram tiap tahun,
roket foton tetap menjadi hitungan di atas kertas, bukan rekayasa.
\end{solution}

\begin{solution}{pb:b3:relativistic-dynamics:1}
\textbf{1.} $E^2 - p^2c^2 = \gamma^2m^2c^4(1 - \beta^2) = m^2c^4$.
\textbf{2.} $E_{\text{tot}}$ dan $\vect p_{\text{tot}}$ bertransformasi
seperti $E$ dan $\vect p$ satu partikel (keduanya jumlahan), sehingga
aljabar yang sama memberikan sebuah invarian; di tempat
$\vect p_{\text{tot}} = \vect 0$ ia menjadi kuadrat energi totalnya:
$Mc^2 = E_{\text{CM}}$.
\textbf{3.} $(E + m_{\text{p}}c^2)^2 - p^2c^2 = m_{\text{p}}^2c^4 +
2Em_{\text{p}}c^2 + m_{\text{p}}^2c^4$.
\textbf{4.} $E \to m_{\text{p}}c^2$ memberikan $M = 2m_{\text{p}}$;
$E \gg m_{\text{p}}c^2$ memberikan $Mc^2 \to \sqrt{2Em_{\text{p}}c^2}$.
\textbf{5.} Gerak nisbi apa pun di antara hasilnya menambahkan energi
kinetik dalam kerangka pusat momentum di atas massa diamnya; $Mc^2$
minimum --- yakni ambangnya --- membuat semuanya diam satu terhadap yang
lain.
\textbf{6.} Momentum laboratorium $\vect p_{\text{tot}} \neq \vect 0$
harus bertahan melewati tumbukannya: hasilnya terhukum untuk bergerak,
dan energi kinetiknya terkunci dari pembuatan massa.
\textbf{7.} $p + p \to p + p + p + \bar p$: muatannya $+2 \to +2$;
bilangan barionnya $+2 \to 1 + 1 + 1 - 1 = +2$. Membuat $\bar p$ sendiri,
atau dengan apa pun yang kurang dari satu barion tambahan penuh,
melanggar salah satu dari keduanya.
\textbf{8.} Empat massa proton yang saling diam: $Mc^2 =
4m_{\text{p}}c^2$.
\textbf{9.} $16m_{\text{p}}^2c^4 = 2m_{\text{p}}^2c^4 +
2Em_{\text{p}}c^2$: $E = 7m_{\text{p}}c^2$, $T = 6m_{\text{p}}c^2 =
\qty{5.63}{GeV}$.
\textbf{10.} Massa diam barunya: $2m_{\text{p}}c^2$ dari $T =
6m_{\text{p}}c^2$ --- sepertiganya; dua pertiga lainnya terbang terus
sebagai energi kinetik hasilnya, yang dilindungi kekekalan momentum.
\textbf{11.} Secara berhadapan, $Mc^2 = 2(T' + m_{\text{p}}c^2) =
4m_{\text{p}}c^2$: $T' = m_{\text{p}}c^2 = \qty{0.94}{GeV}$ tiap berkas
--- enam kali lebih kecil daripada berkas sasaran diam, dan dua belas
kali lebih kecil dalam energi kinetik totalnya.
\textbf{12.} Berkas pada masa itu jauh terlalu renggang untuk
bertumbukan satu sama lain secara berguna; sasaran padat menyediakan
$10^{22}$ proton tiap sentimeter kubik. Harganya: $\qty{5.6}{GeV}$
alih-alih $2 \times \qty{0.94}{GeV}$.
\textbf{13.} $E = 6.2 + 0.94 = \qty{7.14}{GeV}$; $pc =
\sqrt{7.14^2 - 0.938^2} = \qty{7.08}{GeV}$.
\textbf{14.} $p = \num{7.08e9} \times \num{1.6e-19}/\num{3e8} =
\qty{3.8e-18}{kg.m/s}$: $r = p/qB = \qty{15.1}{m}$ --- persis mesin
sebagaimana dibangun.
\textbf{15.} $\beta = pc/E = 0.991$; orbitnya $2\pi r = \qty{95}{m}$:
$f = \qty{3.1}{MHz}$.
\textbf{16.} Pada \qty{10}{MeV}, $\beta = 0.145$: $f =
\qty{0.46}{MHz}$. Ketika energinya memanjat, lajunya (sehingga $f$) dan
momentumnya (sehingga medan yang diperlukan pada jejari tetap)
sama-sama berubah: medan dan frekuensi radionya harus naik bersama ---
yakni muslihat yang mendefinisikan sinkrotron.
\textbf{17.} $\num{6.19e9}/1500 \approx \num{4e6}$ putaran; pada
frekuensi rata-rata berskala megahertz, satu daur pemercepatannya
berorde satu sampai dua sekon.
\textbf{18.} Ambang antiproton $T = 6m_{\text{p}}c^2 =
\qty{5.6}{GeV}$, ditambah kelonggarannya: pembukuan soal ini secara
harfiah adalah alasan pemilihan energi --- dan nama --- mesin itu.
\textbf{19.} $r = p/qB$ tidak memuat massa: magnetnya melenturkan
momentum yang sama secara sama, partikel apa pun itu --- sehingga
berkas terpilihnya tetap mencampur pion, kaon dan (jarang) antiproton,
semuanya pada $\qty{1.19}{GeV}/c$.
\textbf{20.} $\bar p$: $E = \sqrt{1.19^2 + 0.938^2} =
\qty{1.51}{GeV}$, $\beta = 0.79$. $\pi^-$: $E = \sqrt{1.19^2 +
0.140^2} = \qty{1.20}{GeV}$, $\beta = 0.993$.
\textbf{21.} $t_{\bar p} = 12/(0.785c) = \qty{51}{ns}$; $t_\pi =
\qty{40}{ns}$: selisih \qty{11}{ns} menuntut pewaktuan berskala
nanosekon --- yang tersedia, tepat pas, pada 1955.
\textbf{22.} Pencacahnya diatur agar menyala bagi pion yang cepat dan
tetap gelap bagi antiproton yang lambat (veto kecepatan); dengan
perbandingan $1$ berbanding $44000$, hanya dua sidik yang saling bebas (waktu terbang
\emph{dan} ambang Cerenkov) yang dapat menyingkirkan kebetulan palsu.
\textbf{23.} Tepat di atas ambangnya hasilnya lahir hampir dalam
keadaan saling diam: reaksinya nyaris tak mempunyai \omterm{def:b3:hamiltonian-mechanics:hamiltonian}{ruang fase},
sedangkan produksi pion, yang jauh di atas ambang\emph{nya} sendiri,
melimpah --- kelangkaannya dijamin oleh kinematika yang sama yang
menetapkan harganya.
\textbf{24.} $2m_{\text{p}}c^2 = \qty{1.9}{GeV}$ tiap pemusnahan,
lazimnya menjadi segenggam pion yang meluruh lebih lanjut menjadi foton,
muon dan neutrino.
\textbf{25.} Kekekalan \omterm{def:b3:relativistic-dynamics:four-momentum}{momentum-empat} menetapkan harga antiproton pada
$6m_{\text{p}}c^2 = \qty{5.6}{GeV}$ di sasaran diam; Bevatron
\qty{6.2}{GeV} berjejari \qty{15.1}{m} membayarnya; waktu terbang dan
veto Cerenkov menemukan satu partikel antidunia tiap $44000$ penipu
--- dan Hadiah Nobel 1959 mengesahkan pembukuannya.
\end{solution}
